\HETableTitle{3.1e}{Critical risks for implementation}
\label{tab:risks}
\noindent
\begin{tblr}{
  width=\linewidth,
  colspec={X[2,l] X[0.9,l] X[2,l]},
  hlines, vlines,
  row{1}={font=\bfseries,bg=HETableGrey},
}
{Description of risk; indicate (i) likelihood and (ii) severity:
Low / Medium / High} & Work package(s) involved & Proposed risk-mitigation measures \\
\HEPlaceholder{risk; likelihood; severity} & \HEPlaceholder{WP(s)} &
\HEPlaceholder{mitigation measures} \\
\HEPlaceholder{risk; likelihood; severity} & \HEPlaceholder{WP(s)} &
\HEPlaceholder{mitigation measures} \\
\end{tblr}
\begin{heguidance}
A critical risk is a plausible event or issue with a potentially high adverse
impact on the project's ability to achieve its objectives. Likelihood is the
estimated probability of occurrence after accounting for the mitigating
measures. Severity expresses the relative seriousness and significance of the
effect. Rate each as low, medium or high. The list and mitigation measures can
be updated as the project progresses.
\end{heguidance}

